\documentclass[11pt]{article}

\usepackage[margin=1in]{geometry}
\usepackage{amsmath,amssymb,mathtools}
\usepackage{booktabs}
\usepackage{xcolor}
\usepackage[hidelinks]{hyperref}

\newcommand{\E}{\mathbb{E}}
\newcommand{\Normal}{\mathcal{N}}
\newcommand{\Lflow}{\mathcal{L}_{\mathrm{flow}}}

\title{How the MTQL Flow Teacher Is Trained}
\author{}
\date{}

\begin{document}
\maketitle

\section{The short answer}

The action chunk sampled from the replay buffer is the clean endpoint of a path.
It is \emph{not} directly used as the teacher's output label.  The code first
constructs a straight path from the data action chunk to Gaussian noise, and
then uses the derivative of that path as the target velocity:

\begin{equation}
    \boxed{
    x_t=(1-t)a+t\epsilon
    \quad\Longrightarrow\quad
    u_t=\frac{d x_t}{dt}=\epsilon-a
    }
    \label{eq:key}
\end{equation}

Thus, the teacher's supervised training pair is

\begin{equation}
    \boxed{
    (o,h,x_t,t) \longmapsto u_t=\epsilon-a
    }.
\end{equation}

Here, $o$ is the current observation, $h$ is its history, and $a$ is the
corresponding action chunk stored in the replay buffer.

\section{Construction of one training example}

Let the action chunk have $K$ time steps and action dimension $d$:

\begin{equation}
    a\in\mathbb{R}^{K\times d}.
\end{equation}

For every item in a sampled minibatch, the code performs the following steps.

\subsection*{Step 1: take a data action chunk}

\begin{equation}
    a \leftarrow \texttt{batch['actions']}.
\end{equation}

This is an action chunk stored in the offline or online replay buffer.

\subsection*{Step 2: sample a noise endpoint}

\begin{equation}
    \epsilon\sim\Normal(0,I),
    \qquad
    \epsilon\in\mathbb{R}^{K\times d}.
\end{equation}

The noise has exactly the same shape as the action chunk.

\subsection*{Step 3: sample a position along the path}

The implementation samples one scalar $t$ per batch item:

\begin{equation}
    t=0.999\,\widetilde{t}+0.001,
    \qquad
    \widetilde{t}\sim\operatorname{Beta}(1.5,1).
\end{equation}

That scalar is broadcast across all $K\times d$ entries of the action chunk.

\subsection*{Step 4: construct the partially noised input}

\begin{equation}
    x_t=(1-t)a+t\epsilon.
    \label{eq:path}
\end{equation}

This is a convex interpolation, not the sum $a+\epsilon$:

\begin{align}
    t=0 &\quad\Longrightarrow\quad x_0=a,
    \\
    0<t<1 &\quad\Longrightarrow\quad x_t
        \text{ is partly action and partly noise},
    \\
    t=1 &\quad\Longrightarrow\quad x_1=\epsilon.
\end{align}

\subsection*{Step 5: differentiate the chosen path}

This is the crucial construction.  Differentiating
Equation~\eqref{eq:path} with respect to $t$ gives

\begin{align}
    \frac{d x_t}{dt}
    &=\frac{d}{dt}\left((1-t)a+t\epsilon\right)
    \\
    &=-a+\epsilon
    \\
    &=\epsilon-a.
\end{align}

Therefore the exact velocity of the chosen straight path is

\begin{equation}
    \boxed{u_t=\epsilon-a.}
\end{equation}

The target is not guessed, annotated, or obtained from a separate model.  It
is known analytically because the training procedure itself selected both
endpoints $a$ and $\epsilon$ and defined the path between them.

\section{What the teacher predicts}

The flow teacher receives the observation context, the interpolated action
chunk, and the time:

\begin{equation}
    \widehat{v}_t
    =v_{\theta_T}(o,h,x_t,t).
\end{equation}

Its output has the same shape as $a$, $\epsilon$, $x_t$, and $u_t$:

\begin{equation}
    \widehat{v}_t\in\mathbb{R}^{K\times d}.
\end{equation}

The flow-matching objective is

\begin{equation}
    \boxed{
    \Lflow
    =\E_{(o,h,a),\epsilon,t}
      \left[
      \left\|
      v_{\theta_T}(o,h,x_t,t)-(\epsilon-a)
      \right\|_2^2
      \right].
    }
\end{equation}

Across many data actions, noise samples, and times, this regression teaches
the network the conditional vector field associated with the action
distribution given $(o,h)$.  No reward or $Q$ value appears in this loss.

\section{Numerical example}

Consider a two-dimensional action rather than a full action chunk:

\begin{equation}
    a=egin{bmatrix}2\\1\end{bmatrix},
    \qquad
    \epsilon=egin{bmatrix}-1\\3\end{bmatrix},
    \qquad
    t=0.25.
\end{equation}

The input shown to the teacher is

\begin{align}
    x_{0.25}
    &=0.75a+0.25\epsilon
    \\
    &=0.75\begin{bmatrix}2\\1\end{bmatrix}
      +0.25\begin{bmatrix}-1\\3\end{bmatrix}
    \\
    &=\begin{bmatrix}1.25\\1.50\end{bmatrix}.
\end{align}

The regression target is the path velocity

\begin{equation}
    u_t=\epsilon-a
    =\begin{bmatrix}-1\\3\end{bmatrix}
     -\begin{bmatrix}2\\1\end{bmatrix}
    =\begin{bmatrix}-3\\2\end{bmatrix}.
\end{equation}

So this single supervised example is conceptually

\begin{equation}
    \left(o,h,
    x_t=\begin{bmatrix}1.25\\1.50\end{bmatrix},
    t=0.25\right)
    \longmapsto
    \begin{bmatrix}-3\\2\end{bmatrix}.
\end{equation}

\section{Why generation runs from noise to actions}

The velocity $\epsilon-a$ is defined in the forward direction from data at
$t=0$ to noise at $t=1$.  Generation starts at the opposite endpoint,

\begin{equation}
    x_1=\epsilon,
\end{equation}

and integrates using a negative time step.  If $\Delta t>0$, one reverse Euler
step is

\begin{equation}
    x_{t-\Delta t}
    =x_t-\Delta t\,
      v_{\theta_T}(o,h,x_t,t).
\end{equation}

After repeatedly stepping from $t=1$ to $t=0$, the final $x_0$ is the generated
action chunk.  Different initial noise chunks can produce different action
chunks for the same observation context.

\section{Summary of the four quantities}

\begin{center}
\begin{tabular}{lll}
\toprule
Quantity & Meaning & Role \\
\midrule
$a$ & replay action chunk & clean endpoint \\
$\epsilon$ & sampled Gaussian chunk & noise endpoint \\
$x_t=(1-t)a+t\epsilon$ & interpolated chunk & teacher input \\
$u_t=\epsilon-a$ & derivative of the path & teacher label \\
\bottomrule
\end{tabular}
\end{center}

\section{Code correspondence}

The construction appears in
\texttt{agents/mtql\_transformer.py}, lines 138--159:

\begin{verbatim}
actions = batch['actions']
noise = jax.random.normal(noise_rng, actions.shape)
x_t = time_expanded * noise + (1 - time_expanded) * actions
u_t = noise - actions
pred_vel = self.network.select('actor_bc_flow')(..., x_t, time[:, None], ...)
loss = jnp.mean((pred_vel - u_t) ** 2)
\end{verbatim}

\begin{center}
\fbox{%
\parbox{0.88\linewidth}{
The replay action $a$ supplies the clean endpoint.  The noised input is the
interpolation $x_t=(1-t)a+t\epsilon$.  The target is the analytically known
path derivative $u_t=\epsilon-a$.
}}
\end{center}

\end{document}
